\section*{Resumen ejecutivo}

\textbf{El informe distingue dos medidas de la productividad laboral.} Con los datos del DANE para Colombia, descompone el crecimiento del valor agregado real por hora trabajada y por persona empleada entre 2005 y 2025. Las dos medidas pueden moverse en direcciones distintas cuando cambian las horas trabajadas por persona.

\textbf{En el conjunto del periodo, el mayor aporte provino de los servicios de capital.} La productividad creció en promedio 1,03\% anual por hora y 0,88\% por persona, en tasas logarítmicas. El capital aportó 0,79 y 0,73 puntos porcentuales anuales, respectivamente. La composición del trabajo aportó 0,37 puntos a cada medida y la productividad total de los factores (PTF) restó 0,14 puntos.

\textbf{Desde 2022 aumentó el valor agregado por hora y cayó por persona.} Entre 2022 y 2025 las tasas promedio fueron 0,31\% por hora y $-1{,}17$\% por persona. La brecha corresponde a una caída de 1,48\% anual en las horas por persona. En la descomposición por persona, esa caída se reparte entre el aporte directo de las horas y la diferencia entre el capital por persona y el capital por hora.

\textbf{La diferencia se mantuvo en 2025.} El valor agregado por hora creció 0,40\% y el valor agregado por persona cayó 0,56\%. La PTF aportó 0,75 puntos a ambas medidas. En la medición por persona, los aportes negativos de las horas, el capital y la composición del trabajo superaron esa contribución.

\textbf{La medida pertinente depende de la comparación.} El valor agregado por hora indica cuánto se produce en promedio por unidad de tiempo trabajado. El valor agregado por persona indica cuánto se produce en promedio por cada persona empleada durante el año. La descomposición permite identificar qué componentes acompañaron los cambios de cada medida, sin establecer por qué cambiaron las horas ni cuánto de la PTF provino de avances tecnológicos.
\clearpage
